% !Mode:: "TeX:UTF-8"
% VL 15, gebaut aus part-architekturen-anwendungen/25-destillation.md und dem
% Vibecoding-Abschnitt aus 24-anwendungsbeispiele.md
% (Dialogmuster und Zeiten: transcripts/vl14-destillation-2025/notiz.md).
% !Mode:: "TeX:UTF-8"
\newcommand{\thepage}{} % to silence a hyperref warning about pdfpagelabels
\RequirePackage[]{silence}
\WarningsOff[hyperref]
\PassOptionsToPackage{log-declarations=false}{xparse}
%\RequirePackage[l2tabu, orthodox]{nag}
\RequirePackage{luatex85} % to make \pdf...= commands work with LuaTeX
\documentclass[aspectratio=169,t]{beamer}
\usepackage{ifthen}
\usepackage[T1]{fontenc}
\usepackage{iftex}
\ifLuaTeX
 \usepackage{fontspec}
 \setmainfont[Ligatures=TeX]{Source Sans Pro}
\else
 \usepackage[utf8]{inputenc}
\fi
\usepackage{german}
\usepackage{newunicodechar}
\newunicodechar{€}{\texteuro}
\newunicodechar{ß}{\ss}
\renewcommand{\SS}{\iffontchar\font"1E9E \symbol{"1E9E}\else SS\fi}
\usepackage{ellipsis}
\usepackage{microtype}
\usepackage{listings}
\usepackage{sourcecodepro}
\lstset{language=Python, basicstyle=\ttfamily}
\usepackage{amsfonts}
\usepackage{amssymb}
\usepackage{amsthm}
\usepackage{fancyhdr}
\usepackage{graphicx}

\newtheorem{proposition}[theorem]{Proposition}

\DeclareMathOperator*{\argmax}{argmax}
\DeclareMathOperator{\supp}{supp}
\newcommand{\E}{\ensuremath{\mathbb{E}}}
\newcommand{\V}{\ensuremath{\mathbb{V}}}

\setbeamercolor{emph}{fg=purple}
\renewcommand<>{\emph}[1]{%
  {\usebeamercolor[fg]{emph}\only#2{\itshape}#1}%
}



\hypersetup{draft=false,colorlinks=false,pdfpagemode=UseOutlines,plainpages=false,
            breaklinks=true,bookmarksnumbered=true,pdfborder=0 0 0,
            pdftitle={Einführung in Deep Learning WiSe 2026/27 HHU},
            pdfauthor={Konrad Voelkel},
            pdfsubject={data science}}

%\documentclass[aspectratio=169,t]{beamer}
\usetheme{hhubeamer} % put the hhu-beamer-theme in /usr/share/texmf/tex/latex/beamer/themes/

% -------------------- Code listings config --------------------
% Define custom colors for syntax highlighting
\definecolor{codegreen}{rgb}{0,0.6,0}
\definecolor{codegray}{rgb}{0.5,0.5,0.5}
\definecolor{codepurple}{rgb}{0.58,0,0.82}
\definecolor{backcolour}{rgb}{0.96,0.96,0.94} % Very light gray/beige
% Configure the listings package
\lstset{
    language=Python,
    backgroundcolor=\color{backcolour},   
    commentstyle=\color{codegreen}\itshape,
    keywordstyle=\color{blue}\bfseries,
    numberstyle=\tiny\color{codegray},
    stringstyle=\color{codepurple},
    basicstyle=\ttfamily\footnotesize, % Use a monospaced font
    breakatwhitespace=false,         
    breaklines=true,      
             % Enable automatic line breaking
    captionpos=b,                      % Caption at the bottom
    keepspaces=true,                 
    numbers=left,                      % Line numbers on the left
    numbersep=10pt,                    % Distance of line numbers from the code
    showspaces=false,                
    showstringspaces=false,            % Don't show weird underscores in strings
    showtabs=false,                  
    tabsize=4,
    frame=single,                      % Add a subtle border around the code
    rulecolor=\color{black!20},        % Border color
  }

\event[]{}
\author[Wintersemester 2026/27]{Dr. Konrad Völkel, Wintersemester 2026/27}

%\title{Data Science, VL 1}
%\date{16.04.2026}


%%% Local Variables:
%%% mode: latex
%%% TeX-master: t
%%% End:


% Fassung mit Dozentennotizen (zweiter Schirm): die nächste Zeile einkommentieren.
% \setbeameroption{show notes on second screen=right}

\graphicspath{{../images/}}

\title[Destillation (VL 15)]{Destillation -- und Vibecoding\\\large Einführung in Deep Learning, VL 15}
\date{03.02.2027}

\begin{document}

\begin{frame}
  \titlepage

  \note{\textbf{Ablauf, 90 Minuten} (2025: 51 Min. Folien, 7 Min. Fragerunden,
    25 Min. Vibecoding live):\\[2pt]
    \begin{tabular}{@{}rll@{}}
      0--4   & Folie 2       & Letzte Vorlesung, Klausurhinweise\\
      4--18  & Folien 3--5   & \textbf{Dialog: Was lernt ein Netz eigentlich?}\\
      18--28 & Folien 6--7   & Suche gegen Regression, Model Compression\\
      28--38 & Folien 8--9   & Pruning, Destillation: drei Ansätze\\
      38--52 & Folien 10--12 & Distillation Loss, Dark Knowledge\\
      52--60 & Folien 13--14 & DistilBERT, Schüler schlägt Lehrer\\
      60--85 & Folie 15      & \textbf{Vibecoding live}\\
      85--90 & Folie 16      & Klausur, Abschied\\
    \end{tabular}\\[4pt]
    \textbf{Vorbereitung für den Vibecoding-Teil:} Gemini-Konto, der Katzen-Hunde-Datensatz
    (tausend Bilder) und ein vorab erzeugtes, \emph{gelaufenes} Skript. 2025 ging die
    zweite Runde live schief -- das war lehrreich, kostet aber Zeit.}
\end{frame}

\begin{frame}
  \frametitle{Letzte Vorlesung}

  \begin{center}
    \Large
    Klausur: Do 11.02.2027, 9:00--10:30\\[10pt]
    Hörsäle 5F und 2D, Sitzplatzzuteilung vorab
  \end{center}

  \note{Die Probeklausur ist besprochen. Heute noch ein Thema, das erklärt, warum die
    kleinen Modelle von heute so gut sind -- und danach ein Blick darauf, wie man mit
    generiertem Code arbeitet.}
\end{frame}

\begin{frame}
  \frametitle{Was lernt ein Netz eigentlich?}

  \begin{center}
    \Huge Eine Funktion.
  \end{center}

  \note{\textbf{Dialogeinstieg.} Die Frage in den Raum stellen und die Antworten sortieren
    -- 2025 kamen „Datenrepräsentation“ und „Interpretation“, und daraus lassen sich die
    beiden Ansätze der nächsten Folien bauen.}
\end{frame}

\begin{frame}
  \frametitle{Ansatz 1: die Funktion ist bekannt}

  \begin{center}
    \Large
    $y_i = f(x_i)$ liegt vor oder ist billig zu berechnen\\[16pt]
    \normalsize $\Rightarrow$ gewöhnliche Regression
  \end{center}

  \note{Das ist der einfache Fall: Wir kennen die Zielfunktion und approximieren sie.
    Die universelle Approximation aus VL 7 sagt, dass das geht.}
\end{frame}

\begin{frame}
  \frametitle{Ansatz 2: die Funktion ist unbekannt}

  \begin{center}
    \Large
    $f(x_1,\dots,x_k) = P(x_k \mid x_1,\dots,x_{k-1})$
  \end{center}

  \bigskip
  \begin{center}
    \normalsize Das Netz muss die Funktion erst \emph{finden}.
  \end{center}

  \note{Next token prediction -- man könnte die Wahrscheinlichkeit aus Daten schätzen,
    aber das ist teuer.
    Der Autoencoder ist derselbe Fall: Wir trainieren auf die Identität, aber was wir
    eigentlich wollen, ist der Encoder in der Mitte -- eine Funktion, die wir vorher
    nicht kannten.}
\end{frame}

\begin{frame}
  \frametitle{Der rote Faden: Suche gegen Regression}

  \begin{center}
    \Large
    \emph{Suchen} braucht Platz und Flexibilität\\
    $\Rightarrow$ große, tiefe Netze\\[16pt]
    \emph{Regression} auf etwas Bekanntes braucht wenig\\
    $\Rightarrow$ kleine, flache Netze genügen
  \end{center}

  \note{\textbf{Das ist der Satz, an dem die ganze Sitzung hängt.}
    Bild aus 2025, das gut trägt: ein Gebirge mit Zäunen -- je mehr Platz man hat, desto
    leichter findet man einen Weg hindurch. Deshalb lernen große Netze leichter.
    Und daraus folgt die Idee: erst groß suchen, dann klein nachbauen.}
\end{frame}

\begin{frame}
  \frametitle{Model Compression}

  \begin{center}
    \Large
    \textbf{Pruning}: unwichtige Teile abschneiden\\[14pt]
    \textbf{Destillation}: ein kleines Netz vom großen lernen lassen
  \end{center}

  \note{Pruning ist schwierig: Es ist unklar, was weg kann, und nach dem Entfernen werden
    „unwichtige“ Teile manchmal wichtig.
    Historischer Vorläufer: Buciluă, Caruana, Niculescu-Mizil, \emph{Model Compression}
    (2006) -- großes Ensemble trainieren, synthetische Daten labeln lassen, kleines Netz
    darauf trainieren. Das Training wirkt regularisierend.}
\end{frame}

\begin{frame}
  \frametitle{Destillation: drei Ansätze}

  \begin{enumerate}\setlength{\itemsep}{8pt}
  \item Lehrer-Labels auf echten Daten
  \item Synthetische Daten vom Lehrer
  \item \emph{Logits} des Lehrers -- Destillation im engeren Sinn
  \end{enumerate}

  \note{Der dritte ist der interessante. Die ersten beiden sind „der Lehrer labelt“, der
    dritte ist „der Schüler hört dem Lehrer beim \emph{Denken} zu“.
    Hinton, Vinyals, Dean: \emph{Distilling the Knowledge in a Neural Network} (2015).}
\end{frame}

\begin{frame}
  \frametitle{Distillation Loss}

  \begin{center}
    \Large
    $q_i = \dfrac{\exp(z_i/T)}{\sum_j \exp(z_j/T)}$\\[16pt]
    $L_{\text{distill}} = T^2 \cdot D_{KL}\bigl(q^{\text{Lehrer}} \,\|\, q^{\text{Schüler}}\bigr)$
  \end{center}

  \note{Temperature Scaling wie in VL 14 -- mit $T > 1$ werden die kleinen
    Wahrscheinlichkeiten der falschen Klassen sichtbar.
    Der Faktor $T^2$ gleicht aus, dass die Gradienten mit $1/T^2$ skalieren; sonst müsste
    man die Lernrate mit der Temperatur nachziehen.
    In der Praxis kombiniert man das mit gewöhnlicher Kreuzentropie auf den echten Labels.}
\end{frame}

\begin{frame}
  \frametitle{Dark Knowledge}

  \begin{center}
    \Large
    Ein Hund, falsch als „Katze“ gelabelt.\\[16pt]
    Der Lehrer sagt „Katze“ -- aber \emph{vor} dem argmax\\
    steht $P(\text{Hund})$ deutlich über null.
  \end{center}

  \note{\textbf{Das ist die Take-Home-Message der Sitzung:} Der Schüler hört dem Lehrer
    nicht beim \emph{Urteilen} zu, sondern beim \emph{Denken}. Die weichen
    Wahrscheinlichkeiten vor dem argmax tragen die Information.
    Deshalb kann ein Schüler auch mit fehlerhaften Labels gut werden -- und deshalb
    braucht man die Temperatur.}
\end{frame}

\begin{frame}
  \frametitle{Der Schüler kann den Lehrer übertreffen}

  \begin{center}
    \Large
    \emph{Born Again Neural Networks} (2018)
  \end{center}

  \bigskip
  \begin{center}
    \normalsize sogar bei \emph{gleicher} Architektur.
  \end{center}

  \note{Das ist die Frage, die 2025 aus dem Hörsaal kam: Darf der Schüler dem Lehrer
    widersprechen? Antwort: Er wird sogar besser, wenn er es tut -- die weichen Label
    wirken wie eine Regularisierung.
    Eine zweite Frage von 2025: Warum gibt es dann überhaupt noch große Modelle? Weil
    jemand zuerst \emph{suchen} muss.}
\end{frame}

\begin{frame}
  \frametitle{DistilBERT}

  \begin{center}
    \Large
    40\,\% kleiner, 97\,\% der Leistung\\[14pt]
    \normalsize jeder zweite Layer als Initialisierung,\\
    drei Loss-Terme
  \end{center}

  \note{Distillation Loss (KL zu den skalierten Logits), Masked Language Modeling Loss wie
    bei BERT, und ein Cosine Loss, der die Richtung der hidden states angleicht.
    Das ist das Modell aus dem Task-Rundgang in VL 3 -- der Kreis schließt sich.
    Weitere Beispiele, falls Zeit: Speculative Decoding, SDXL Turbo, Gemini Flash,
    Chain of Thought als Destillationsziel.}
\end{frame}

\begin{frame}
  \frametitle{Vibecoding}

  \begin{center}
    \Large
    Aufgabe: Katzen und Hunde sortieren -- mit generiertem Code.
  \end{center}

  \note{\textbf{Live-Teil, etwa 25 Minuten.} Vorbereitet: Prompt und erstes Skript, schon
    einmal gelaufen. Im Editor in \emph{Ausführungsreihenfolge} lesen, nicht von oben nach
    unten; Ergebnis im Dateimanager prüfen.
    \textbf{Die Pointe, die stehen bleiben muss:} Hohe Accuracy sagt nichts, wenn die
    Labelphase Murks war. Generierter Code läuft fast immer -- ob er das Richtige tut,
    sieht man nur, wenn man die Annahmen im Code kennt und in die Ergebnisordner guckt.
    2025 ist die zweite Runde live schiefgegangen (vermutlich nicht gelöschte Kopien im
    sortierten Ordner). Das war lehrreich -- aber besser eingeplant als überraschend.}
\end{frame}

\begin{frame}
  \frametitle{Das war's}

  \begin{center}
    \Large
    Klausur: Do 11.02.2027, 9:00--10:30\\[14pt]
    \normalsize Fragen bis dahin: im Chat oder in der Fragestunde.
  \end{center}

  \bigskip
  \begin{center}
    Viel Erfolg!
  \end{center}

  \note{Kurz zusammenfassen, was das Semester getragen hat: Skalarprodukt plus
    Nichtlinearität, eine Loss-Funktion minimieren, Backpropagation als Buchhaltung,
    latente Räume als das, worin alles lebt.
    Und ehrlich sagen, was wir nicht gemacht haben: Zeitreihen, Reinforcement Learning,
    eigene Transformer trainieren -- dafür gibt es die Aufbaumodule im Master.}
\end{frame}

\end{document}

%%% Local Variables:
%%% mode: latex
%%% TeX-master: t
%%% End:
